\section{排队系统与 M/M/1 模型}
\label{sec:06-Random-Processes/05-Continuous-Time-Markov-Chains-and-Queues/04}

一个服务的 CPU 利用率从 \(70\%\) 涨到 \(85\%\)，平均延迟翻了一倍。你提扩容，得到的答复是：还有 \(15\%\) 的余量，先优化代码。

这个答复背后有一个隐含的线性假设——利用率还没到顶，延迟就不该出事。这一节要说明那个假设错在哪里，错得有多离谱，以及为什么容量规划不能按平均负载来做。

顺带还会拿到一件工具，它比排队模型本身管用得多：一条几乎不需要任何假设的等式，把吞吐、并发和延迟三个数锁在一起。

\subsection{Little 定律：三个数只有两个是自由的}

先不谈任何分布。设想任意一个有进有出的系统——一段服务、一个线程池、一整条微服务链路，甚至一家餐馆。记 \(L\) 为长期平均在系统中的顾客数，\(\lambda\) 为长期平均到达率，\(W\) 为每位顾客的平均逗留时间。

\begin{theorem}
若系统长期稳定（到达率、逗留时间的长期平均都存在且有限，且没有顾客永远滞留），则
\[
L=\lambda W.
\]
\end{theorem}

条件里没有出现任何分布。到达可以不是 Poisson，服务时间可以任意，可以有优先级、可以抢占、可以多台服务器、可以是一整个由几十个服务组成的调用图——只要它稳定，等式就成立。这种近乎无条件的普适性，来自证明只用了一件事：同一块面积可以按两个方向来数。

\begin{center}
\begin{tikzpicture}[x=1cm,y=1cm,font=\small,>=stealth]
  \fill[blue!14] (0.4,0) rectangle (1.0,0.55);
  \fill[blue!14] (1.0,0) rectangle (1.8,1.1);
  \fill[blue!14] (1.8,0) rectangle (2.3,0.55);
  \fill[blue!14] (2.3,0) rectangle (2.8,1.1);
  \fill[blue!14] (2.8,0) rectangle (3.4,1.65);
  \fill[blue!14] (3.4,0) rectangle (4.0,1.1);
  \fill[blue!14] (4.0,0) rectangle (4.6,0.55);
  \draw[blue!75,thick] (0,0) -- (0.4,0) -- (0.4,0.55) -- (1.0,0.55) -- (1.0,1.1)
    -- (1.8,1.1) -- (1.8,0.55) -- (2.3,0.55) -- (2.3,1.1) -- (2.8,1.1)
    -- (2.8,1.65) -- (3.4,1.65) -- (3.4,1.1) -- (4.0,1.1) -- (4.0,0.55)
    -- (4.6,0.55) -- (4.6,0) -- (5.2,0);
  \draw[->,black!55] (0,0) -- (5.4,0) node[right,font=\footnotesize] {\(t\)};
  \draw[->,black!55] (0,0) -- (0,2.2) node[above,font=\footnotesize] {\(N(t)\)};
  \node[font=\footnotesize,align=center,black!80] at (2.6,-0.7) {横着数：\\面积 \(=\) 平均人数 \(\times\) 时长};
  \draw[black!25] (6.0,-1.1) -- (6.0,2.4);
  \begin{scope}[shift={(6.8,0)}]
    \fill[blue!14] (0.4,0.25) rectangle (1.8,0.6);
    \fill[blue!14] (1.0,0.7) rectangle (3.4,1.05);
    \fill[blue!14] (2.3,1.15) rectangle (4.0,1.5);
    \fill[blue!14] (2.8,1.6) rectangle (4.6,1.95);
    \draw[blue!75] (0.4,0.25) rectangle (1.8,0.6);
    \draw[blue!75] (1.0,0.7) rectangle (3.4,1.05);
    \draw[blue!75] (2.3,1.15) rectangle (4.0,1.5);
    \draw[blue!75] (2.8,1.6) rectangle (4.6,1.95);
    \draw[->,black!55] (0,0) -- (5.4,0) node[right,font=\footnotesize] {\(t\)};
    \draw[->,black!55] (0,0) -- (0,2.2);
    \foreach \y/\lab in {0.42/1,0.87/2,1.32/3,1.77/4}
      \node[left,font=\footnotesize,black!65] at (-0.08,\y) {\(\lab\)};
    \node[font=\footnotesize,align=center,black!80] at (2.6,-0.7) {竖着数：\\面积 \(=\) 人数 \(\times\) 平均逗留};
  \end{scope}
\end{tikzpicture}

\emph{左边是系统内人数随时间的阶梯，右边是同样这四位顾客各自的逗留区间。两张图涂的是同一块面积——左图把它按时间切成竖条，右图把它按顾客切成横条。面积不在乎你从哪个方向切，这就是 \(L=\lambda W\) 的全部内容。}
\end{center}

骨架写出来只有三行。在 \([0,T]\) 上考虑积分 \(\int_{0}^{T} N(t)\,dt\)。按时间读，它等于 \(T\) 乘以这段时间内的平均人数；按顾客读，每位顾客对积分的贡献恰好是他自己的逗留时长，于是它等于这段时间内的到达人数乘以平均逗留时长。两式相等，除以 \(T\) 并令 \(T\to\infty\)，左边趋于 \(L\)，右边的到达人数除以 \(T\) 趋于 \(\lambda\)，平均逗留趋于 \(W\)。

条件对账也就这么两条。稳定性用在极限存在上——若队列在无限增长，左边的平均人数发散，等式两边都没有意义。边界一致性用在``顾客''的定义上：\(L\)、\(\lambda\)、\(W\) 必须指同一个系统的同一批顾客。这一条在实践中比稳定性更容易翻车。

翻车的样子通常是这样：你从服务端指标读到平均并发 40，从网关读到吞吐每秒 2000，两数一除得到 20 毫秒，然后发现和客户端测的 P50 对不上。原因往往是三个数量的边界不同——并发只统计了进入线程池之后的请求，吞吐把重试算了两次，而客户端的延迟从 DNS 解析就开始计时。等式没错，是三个数说的不是同一个系统。

边界对齐之后，这条式子在日常里极其顺手：吞吐和并发通常是现成的监控指标，延迟反而难测（要打点、要对齐时钟），而 \(W=L/\lambda\) 让你从前两个反推第三个。反过来做容量规划也一样：要把延迟压到 10 毫秒、吞吐要撑住每秒 2000，那么系统里平均只能有 20 个请求在飞——这直接约束了线程池和连接池该开多大。

\subsection{M/M/1 就是最简单的生灭过程}

Little 定律不告诉你 \(W\) 是多少，只告诉你三个数的关系。要预测延迟随负载怎么变，得对到达和服务作出假设。最简单的一组假设写作 M/M/1：三个位置分别是到达过程、服务时间分布、服务台数量；两个 M 都取自 Markovian，意思是 Poisson 到达（等价于指数分布的到达间隔）与独立同分布的指数服务时间；1 是单服务台。

这组假设的作用是让``系统内人数'' \(N(t)\) 自己就构成一个 CTMC。指数分布的无记忆性在这里是关键：不管当前这位顾客已经被服务了多久，剩余服务时间的分布不变，所以只要知道人数就够了。换成一般服务时间，光看人数无法预测下一次离开，还得记住已服务时长，状态就得扩充。

于是 \(N(t)\) 是一个出生率恒为 \(\lambda\)、死亡率恒为 \(\mu\) 的生灭过程（空系统没有服务完成，所以 0 处没有向下的箭头）。上一节的切口论证直接可用：

\[
\pi_n\lambda=\pi_{n+1}\mu
\quad\Longrightarrow\quad
\pi_n=\rho^n\pi_0,
\qquad \rho\eqdef\frac{\lambda}{\mu}.
\]

递推是等比的，归一化就是等比级数求和，收敛当且仅当 \(\rho<1\)。稳定性条件不是额外假设进来的，它是``那串数能不能加成 1''的直接结果。此时 \(\pi_0=1-\rho\)，

\[
\pi_n=(1-\rho)\rho^n.
\]

顺带读出两件事：服务台空闲的概率是 \(1-\rho\)，所以 \(\rho\) 就是利用率；系统内人数服从几何分布，尾巴是指数衰减的。

\(\rho\ge 1\) 时没有平稳分布，但两种失败并不一样。\(\lambda>\mu\) 时队列有明确的漂移方向，长度线性增长；\(\lambda=\mu\) 时链是零常返的——它会无穷多次回到空状态，可平均返回时间是无穷，队长做的是无边界的随机游走。``平均刚好够用''不是勉强及格，它和不够用一样不可用，因为它没给随机波动留任何余量。

\subsection{爆炸曲线}

对几何分布求期望，得到系统内平均人数

\[
L=\sum_{n\ge 0}n(1-\rho)\rho^n=\frac{\rho}{1-\rho}.
\]

这个式子的形状比它的推导重要得多。

\begin{center}
\begin{tikzpicture}[x=1cm,y=1cm,font=\small,>=stealth]
  \draw[->,black!55] (0,0) -- (7.1,0) node[right,font=\footnotesize] {\(\rho\)};
  \draw[->,black!55] (0,0) -- (0,5.1) node[above,font=\footnotesize] {\(L\)};
  \draw[black!60,dashed,thick] (6,0) -- (6,4.9);
  \node[above,font=\footnotesize,black!70] at (6,4.9) {\(\rho=1\)};
  \foreach \r/\lab in {0/0,1.5/0.25,3/0.5,4.5/0.75}
    {\draw[black!45] (\r,0) -- (\r,-0.1);
     \node[below,font=\footnotesize,black!65] at (\r,-0.12) {\(\lab\)};}
  \foreach \l/\lab in {0.45/1,1.8/4,4.05/9}
    {\draw[black!25] (0,\l) -- (6,\l);
     \node[left,font=\footnotesize,black!65] at (-0.08,\l) {\(\lab\)};}
  \draw[blue!75,thick] plot[domain=0:0.905,samples=140] ({6*\x},{0.45*\x/(1-\x)});
  \fill[black!70] (3,0.45) circle (1.8pt);
  \fill[black!70] (4.8,1.8) circle (1.8pt);
  \fill[red!75] (5.4,4.05) circle (2.2pt);
  \node[right,font=\footnotesize,black!75] at (3.1,0.28) {\(\rho=0.5\)};
  \node[right,font=\footnotesize,black!75] at (2.2,1.95) {\(\rho=0.8\)};
  \node[left,font=\footnotesize,red!75] at (5.25,4.3) {\(\rho=0.9\)};
  \draw[->,black!65] (4.85,1.95) -- (5.32,3.88);
  \node[right,font=\footnotesize,black!80,align=left] at (2.0,3.1)
    {利用率 \(+12.5\%\)，\\平均队长 \(\times 2.25\)};
\end{tikzpicture}

\emph{\(L=\rho/(1-\rho)\)。左半段几乎是平的，一路走到 \(\rho=0.5\) 系统里也才平均一个人；越过 \(0.8\) 之后曲线立起来，逼近 \(\rho=1\) 时竖直上天。分母里的 \(1-\rho\) 是剩余余量，队长与余量成反比——余量减半，队长翻倍。容量规划的全部教训在这条曲线的右半段。}
\end{center}

配上 Little 定律，延迟就出来了：

\[
W=\frac{L}{\lambda}=\frac{1}{\mu-\lambda},
\qquad
W_q=W-\frac1\mu=\frac{\rho}{\mu-\lambda}.
\]

第一式的分母同样是余量，不是容量。给个数：服务台每秒能处理 1000 个请求，即纯服务时间 1 毫秒。负载每秒 800 时 \(W=5\) 毫秒，负载 900 时 \(W=10\) 毫秒，负载 950 时 \(W=20\) 毫秒。代码一行没改，服务本身始终是 1 毫秒，端到端延迟却先后是它的 5 倍、10 倍、20 倍——多出来的全是排队。

这就回答了开头那个驳回意见。利用率从 \(0.7\) 走到 \(0.85\)，余量从 \(0.3\) 缩到 \(0.15\)，而 \(W\) 的分母正是余量——延迟不多不少正好翻倍，与观测对得上。同一段路上平均队长 \(L\) 从 \(2.33\) 涨到 \(5.67\)。\(15\%\) 的利用率余量从来不是 \(15\%\) 的性能余量。

\subsection{为什么余量是必需品}

爆炸的机制不难说清。服务能力是长期平均意义上够用的，但请求不会礼貌地均匀到达，服务时间也参差不齐。短时间内的挤兑必然发生，队列就是在这些时刻堆起来的。堆起来的队列靠什么消化？靠后面那些请求稀疏的空隙。而空隙的总量正比于 \(1-\rho\)。

利用率越高，用来还债的空闲时间越少，同样一次突发要花越久才能排干净。到 \(\rho\to 1\)，空闲趋于零，任何一次波动造成的积压都无法在有生之年被消化——这就是分母上那个 \(1-\rho\)。

所以高利用率和低延迟之间存在硬性张力，不是工程实现不够好。规划容量时 \(\lambda<\mu\) 只是入场券；真正要定的是留多少余量，而余量的多少取决于你能接受多大的尾延迟。顺带一提，上面算的全是平均值，尾部比平均更陡：M/M/1 的等待时间尾部按 \(e^{-(\mu-\lambda)w}\) 衰减，衰减率也是那个余量，\(\rho\) 一高，P99 比均值恶化得更快。

\subsection{这个模型什么时候不能用}

M/M/1 的两个 M 都值得怀疑。

服务时间的方差是第一个漏洞。M/G/1 的 Pollaczek--Khinchine 公式说明，平均等待时间不仅取决于 \(\rho\)，还正比于服务时间的二阶矩——同样的平均服务时间，方差越大排队越长。指数分布的变异系数恰好是 1，属于中等；真实服务时间常常更糟，一个偶尔出现的慢查询能把整条链路的延迟拖高。若你的服务时间分布明显重尾，用 M/M/1 估出的延迟会系统性偏乐观。

到达过程是第二个。Poisson 假设意味着到达是完全无记忆的，而真实流量有日内周期、有重试风暴、有上游的批量提交。突发性会让实际队列比模型预测的长。

结构上的偏离还有一串：多服务台该用 M/M/c（并且 \(c\) 台合并成一个队列比 \(c\) 个独立队列好得多），有限缓冲区该用 M/M/1/K（满了就拒绝，此时进入系统的实际到达率不再是 \(\lambda\)，Little 定律里要换成有效到达率），还有优先级、超时放弃、队长相关的服务速率等等。

这些都不妨碍 M/M/1 值得记住。它给的不是精确的延迟预测，而是一个正确的形状：延迟与剩余余量成反比。真实系统只会比它更糟——方差更大、到达更突发、慢请求更慢——所以把它当作乐观的下界来读，判断不会错到哪里去。

至此，状态还是可数的，时间已经连续了。下一章把最后一步也走完：让状态本身连续变化，于是轨道不再是阶梯，而是一条处处抖动、处处不可微的曲线，微分方程需要重新定义。
